% Part: proof-theory
% Chapter: natural-deduction
% Section: translation-G2i

\documentclass[../../../include/open-logic-section]{subfiles}

\begin{document}

\olfileid{pt}{nat}{tgi}

\olsection{ترجمه از \Log{G2i} به \Log{N2i}}

مقدم‌های دنباله‌ها در \Log{G2i} چندمجموعه‌های~$\Gamma$ از
!!{formula}ها هستند، حال آنکه مقدم‌های دنباله‌ها در \Log{N2i}
مجموعه‌های~$\Gamma'$ از فرمول‌های برچسب‌دارند که هر برچسب در آن‌ها فقط
یک بار ظاهر می‌شود. می‌گوییم مجموعهٔ~$\Gamma'$ از !!{formula}های
برچسب‌دار با چندمجموعهٔ~$\Gamma$ \emph{متناظر است} اگر برای هر
!!{formula}~$!A$، شمار !!{element}های~$\Gamma'$ به شکل~$x:!A$ کمتر یا
مساوی چندگانگی~$!A$، یعنی شمار نسخه‌های~$!A$، در~$\Gamma$ باشد.

\begin{prop}\ollabel{prop:G2i-to-N2i}
اگر $\Log{G2i}+\Cut\Proves\Gamma\Sequent\Delta$، آنگاه
$\Log{N2i}\Proves\Gamma'\Sequent !C$، که در آن اگر $\Delta$ تهی باشد
$!C\ident\lfalse$، و اگر $\Delta=\{!A\}$ باشد $!C\ident !A$ است؛ و
$\Gamma'$ با~$\Gamma$ متناظر است.
\end{prop}

\begin{proof}
  نشان می‌دهیم اگر $\pi$ !!a{proof}ی برای $\Gamma\Sequent\Delta$ در
  $\Log{G2i}+\Cut$ باشد، $\Gamma'$ای متناظر با~$\Gamma$ و
  !!a{proof}~$\delta$ای برای~$\Gamma'\Sequent !C$ در~\Log{N2i} وجود
  دارد. بر $n=\pheight{\pi}$ استقرا می‌کنیم.

  اگر $n=0$، $\pi$ یک اصل است. اگر اصل
  $\lfalse\Sequent\quad$ باشد، $\delta$ اصل
  $x:\lfalse\Sequent\lfalse$ است؛ یعنی
  $\Gamma'=\{x:\lfalse\}$ و $!C\ident\lfalse$. اگر $\pi$ اصلی به شکل
  $!A\Sequent !A$ باشد، $\delta$ برابر $x:!A\Sequent !A$ است.

  اگر $n>0$، $\pi$ به قاعده‌ای~$R$ پایان می‌یابد. فرض استقرا وجود
  !!{proof}هایی در~\Log{N2i} را برای دنباله‌های متناظر با مقدمات آن
  قاعده تضمین می‌کند. با تغییر برچسب‌ها و، در صورت لزوم، تغییر نام
  متغیرهای ویژه، فرض می‌کنیم برچسب‌های ترخیص در این برهان‌ها دوبه‌دو
  متمایزند، متغیرهای ویژهٔ نسخه‌های مقدمات نسبت به یکدیگر و نسبت به
  نتیجه تازه‌اند، و برچسب~$x$ فقط بالای استنتاجی که آن را مرخص می‌کند
  ظاهر می‌شود، اگر اصلاً مرخص شود. سپس نشان می‌دهیم چگونه می‌توان این
  !!{proof}ها را برای ساختن برهانی از یک $\Gamma'\Sequent !C$ مناسب
  ترکیب کرد. حالت‌ها را برحسب~$R$ جدا می‌کنیم.

  \begin{enumerate}
    \item اگر $R$ برابر \RightR{\Weakening} باشد، $\pi$ چنین پایان
    می‌یابد:
    \[
    \AxiomC{}
    \RightLabel{$\pi_1$}
    \Deduce$\Gamma \fCenter$
    \RightLabel{\RightR{\Weakening}}
    \UnaryInf$\Gamma \fCenter !A$
    \DisplayProof
    \]
    اعمال فرض استقرا بر~$\pi_1$، !!a{proof}ی برای
    ~$\Gamma'\Sequent\lfalse$ می‌دهد. با اعمال~$\FalseInt$، !!a{proof}ی
    برای $\Gamma'\Sequent !A$ به دست می‌آوریم.

    \item اگر $R$ برابر \LeftR{\Weakening} باشد، $\pi$ چنین پایان
    می‌یابد:
    \[
    \AxiomC{}
    \RightLabel{$\pi_1$}
    \Deduce$\Gamma \fCenter \Delta$
    \RightLabel{\LeftR{\Weakening}}
    \UnaryInf$!A, \Gamma \fCenter \Delta$
    \DisplayProof
    \]
    فرض استقرا بر~$\pi_1$، !!a{proof}ی برای~$\Gamma'\Sequent !C$
    می‌دهد که همان را~$\delta$ می‌گیریم. توجه کنید اگر $\Gamma'$ با
    $\Gamma$ متناظر باشد، با~$!A,\Gamma$ نیز متناظر است: شمار
    $x:!A$ها در~$\Gamma'$ از شمار نسخه‌های~$!A$ در~$\Gamma$ بیشتر
    نیست و بنابراین از شمار نسخه‌های آن در چندمجموعهٔ~$!A,\Gamma$ نیز
    بیشتر نیست.

    \item اگر $R$ برابر \LeftR{\Contraction} باشد، $\pi$ چنین پایان
    می‌یابد:
    \[
    \AxiomC{}
    \RightLabel{$\pi_1$}
    \Deduce$!A, !A, \Gamma \fCenter \Delta$
    \RightLabel{\LeftR{\Contraction}}
    \UnaryInf$!A, \Gamma \fCenter \Delta$
    \DisplayProof
    \]
    فرض استقرا بر~$\pi_1$، !!a{proof}~$\delta_1$ای برای
    $\Pi,\Gamma'\Sequent !C$ می‌دهد که در آن
    $\Pi\subseteq\{x:!A,y:!A\}$. اگر
    $\Pi=\{x:!A,y:!A\}$، پس از تازه‌کردن دیگر برچسب‌ها، همهٔ
    !!{formula}های برچسب‌دار $y:!A$ در~$\delta_1$ را با~$x:!A$
    جایگزین می‌کنیم. چون $x$ در بافت دنبالهٔ پایانی~$\delta_1$ ظاهر
    می‌شود، می‌توان فرض کرد هیچ استنتاجی در~$\delta_1$ !!a{formula}ی
    به شکل~$x:!B$ را مرخص نمی‌کند؛ یعنی هیچ $x:!B$ در سمت چپ دنباله‌ای
    از~$\delta_1$ فعال نیست. در نتیجه، حاصل همچنان !!a{proof}ی در
    ~\Log{N2i} است.

    \item اگر $R$ برابر $\RightR{\lif}$ باشد، $\pi$ چنین پایان می‌یابد:
    \[
    \AxiomC{}
    \RightLabel{$\pi_1$}
    \Deduce$!A, \Gamma \fCenter !B$
    \RightLabel{\RightR{\lif}}
    \UnaryInf$\Gamma \fCenter !A \lif !B$
    \DisplayProof
    \]
    بنا بر فرض استقرا، !!a{proof}~$\delta_1$ای برای
    $x:!A,\Gamma'\Sequent !B$ یا $\Gamma'\Sequent !B$ وجود دارد که
    $\Gamma'$ با~$\Gamma$ متناظر است. برچسب~$x$ را در~$\delta_1$ تازه
    می‌کنیم و~$\delta$ را حاصل اعمال \Intro{\lif} بر~$\delta_1$
    می‌گیریم؛ در حالت نخست~$x$ مرخص می‌شود و در حالت دوم قاعده هیچ
    فرضی را مرخص نمی‌کند.

    \item اگر $R$ برابر \LeftR{\lif} باشد، $\pi$ چنین پایان می‌یابد:
    \[
    \AxiomC{}
    \RightLabel{$\pi_1$}
    \Deduce$\Gamma_1 \fCenter !A$
    \AxiomC{}
    \RightLabel{$\pi_2$}
    \Deduce$!B, \Gamma_2 \fCenter \Delta$
    \RightLabel{\LeftR{\lif}}
    \BinaryInf$!A \lif !B, \Gamma_1, \Gamma_2 \fCenter \Delta$
    \DisplayProof
    \]
    فرض استقرا !!{proof}~$\delta_1$ای برای
    $\Gamma_1'\Sequent !A$ و~$\delta_2$ای برای
    $x:!B,\Gamma_2'\Sequent !C$ یا $\Gamma_2'\Sequent !C$ می‌دهد. اگر
    $\delta_2$ شکل دوم را داشته باشد، خود آن می‌تواند~$\delta$ باشد.
    در غیر این صورت، برچسب‌ها و متغیرهای ویژهٔ~$\delta_1$ و
    ~‌$\delta_2$ را به‌طور سازگار تازه می‌کنیم تا در گرفتن نسخه‌های
    برهان هیچ برچسب یا متغیر ویژه‌ای گرفتار نشود و هیچ برچسبی میان دو
    برهان مشترک نباشد. آنگاه بافتِ پیوستهٔ
    $\Gamma_1',\Gamma_2'$ با چندمجموعهٔ~$\Gamma_1,\Gamma_2$ متناظر
    است. چون $x:!B$ در دنبالهٔ پایانی~$\delta_2$ ظاهر می‌شود، در هیچ
    استنتاج دارای برچسب ترخیص~$x$ در~$\delta_2$ فعال نیست. هر اصل
    $x:!B\Sequent !B$ در~$\delta_2$ را، با نسخه‌ای تازه و
    بی‌گرفتاری، با !!{proof}~$\delta_3$ زیر جایگزین می‌کنیم:
    \[
    \Axiom$x: !A \lif !B \fCenter !A \lif !B$
    \AxiomC{}
    \RightLabel{$\delta_1$}
    \Deduce$\Gamma_1' \fCenter !A$
    \RightLabel{\Elim{\lif}}
    \BinaryInf$x: !A \lif !B, \Gamma_1' \fCenter !B$
    \DisplayProof
    \]
    و هر رخداد $x:!B$ در~$\delta_2$ را با $x:!A\lif !B$ جایگزین
    می‌کنیم. حاصل
    $\Subst{\delta_2}{\delta_3}{x:!B}$، !!a{proof}ی برای
    $x:!A\lif !B,\Gamma_1',\Gamma_2'\Sequent !C$ است.

    \item اگر $R$ برابر \RightR{\land} باشد، $\pi$ چنین پایان می‌یابد:
    \[
    \AxiomC{}
    \RightLabel{$\pi_1$}
    \Deduce$\Gamma_1 \fCenter !A$
    \AxiomC{}
    \RightLabel{$\pi_2$}
    \Deduce$\Gamma_2 \fCenter !B$
    \RightLabel{\RightR{\land}}
    \BinaryInf$\Gamma_1, \Gamma_2 \fCenter !A \land !B$
    \DisplayProof
    \]
    فرض استقرا برهان~$\delta_1$ برای $\Gamma_1'\Sequent !A$ و
    ~‌$\delta_2$ برای $\Gamma_2'\Sequent !B$ می‌دهد. با تازه‌کردن
    برچسب‌ها و متغیرهای ویژه در~$\delta_2$ تضمین می‌کنیم هیچ‌یک با
    ~‌$\delta_1$ مشترک نباشد. آنگاه بافت پیوستهٔ
    $\Gamma_1',\Gamma_2'$ با~$\Gamma_1,\Gamma_2$ متناظر است. با اعمال
    $\Intro{\land}$ بر دنباله‌های پایانی $\delta_1$ و~$\delta_2$،
    !!a{proof}~$\delta$ را به دست می‌آوریم.

    \item اگر $R$ برابر $\LeftR{\land}$ باشد، برای نمونه $\pi$ چنین
    پایان می‌یابد:
    \[
    \AxiomC{}
    \RightLabel{$\pi_1$}
    \Deduce$!A, \Gamma \fCenter \Delta$
    \RightLabel{\LeftR{\land}}
    \UnaryInf$!A \land !B, \Gamma \fCenter \Delta$
    \DisplayProof
    \]
    بنا بر فرض استقرا، !!a{proof}~$\delta_1$ای برای
    $x:!A,\Gamma'\Sequent !C$ یا $\Gamma'\Sequent !C$ وجود دارد و
    $\Gamma'$ با~$\Gamma$ متناظر است. در حالت دوم می‌توانیم
    $\delta=\delta_1$ بگیریم. در حالت نخست، چون $x:!A$ در دنبالهٔ
    پایانی ظاهر می‌شود، در هیچ استنتاج دارای برچسب ترخیص~$x$ در
    ~‌$\delta_1$ فعال نیست. هر اصل $x:!A\Sequent !A$ را با نسخه‌ای
    تازه و بی‌گرفتاری از !!{proof}~$\delta_2$ زیر جایگزین می‌کنیم:
    \[
    \Axiom$x: !A \land !B \fCenter !A \land !B$
    \RightLabel{\Elim{\land}[1]}
    \UnaryInf$x: !A \land !B \fCenter !A$
    \DisplayProof
    \]
    و هر رخداد $x:!A$ را با $x:!A\land !B$ جایگزین می‌کنیم؛ یعنی
    $\Subst{\delta_1}{\delta_2}{x:!A}$ را می‌گیریم. این حاصل
    !!a{proof}ی برای $x:!A\land !B,\Gamma'\Sequent !C$ است.

    \item اگر $R$ برابر \Cut باشد، $\pi$ چنین پایان می‌یابد:
    \[
    \AxiomC{}
    \RightLabel{$\pi_1$}
    \Deduce$\Gamma_1 \fCenter !A$
    \AxiomC{}
    \RightLabel{$\pi_2$}
    \Deduce$!A, \Gamma_2 \fCenter \Delta$
    \RightLabel{\Cut}
    \BinaryInf$\Gamma_1, \Gamma_2 \fCenter \Delta$
    \DisplayProof
    \]
    فرض استقرا برهان~$\delta_1$ برای $\Gamma_1'\Sequent !A$ و
    ~‌$\delta_2$ برای $x:!A,\Gamma_2'\Sequent !C$ یا
    $\Gamma_2'\Sequent !C$ می‌دهد. در حالت دوم، $\delta=\delta_2$
    می‌گیریم. در غیر این صورت، برچسب‌ها و متغیرهای ویژه را در نسخه‌های
    ~‌$\delta_1$ و~$\delta_2$ تازه می‌کنیم و هر اصل
    $x:!A\Sequent !A$ در~$\delta_2$ را به‌گونه‌ای بی‌گرفتاری با نسخه‌ای
    تازه از~$\delta_1$ جایگزین می‌کنیم. در نتیجه
    $\delta=\Subst{\delta_2}{\delta_1}{x:!A}$، !!{proof}ی برای
    $\Gamma_1',\Gamma_2'\Sequent !C$ است.
  \end{enumerate}
  حالت‌های باقی‌مانده مشابه‌اند.
\end{proof}

\end{document}
